\section{Philosophy of Causality}
\label{sec:philosophy}

\subsection{The Classical Philosophy of Causality}
\label{sec:philosophy_foundation}


\subsubsection{Plato}
\label{sec:history_plato}

Plato is believed to be the first to have explored the cause in a systematic way, most notably in his dialogue Timaeus\cite{archer1888timaeus} (c. 360 BC). In this work, Plato explains the creation of the cosmos through the actions of a divine craftsman, the Demiurge. The Demiurge looks to the eternal Forms as a model to bring order to the pre-existing, chaotic matter. Plato’s concept of causality is thus tied to his broader metaphysical theory of Forms, where true causes are the eternal patterns or blueprints of which the physical world is merely a copy. He identifies several contributing factors necessary for the creation of the world, including the Demiurge (the efficient cause), the Forms (the formal cause), and the Receptacle (the material space)\cite{archer1888timaeus}. 

\subsubsection{Aristotle}
\label{sec:history_aristotle}

A student of Plato, Aristotle (c. 350 BC) formalized the notion of causality in his Physics Book II with the ``Four Causes''\cite{falcon2006aristotlecausality}:

\begin{enumerate}
    \item The material cause or that which is given in reply to the question, ``What is it made out of?''
    \item The formal cause or that which is given in reply to the question, ``What is it?''. What is singled out in the answer is the essence of the what-it-is-to-be something.
    \item The efficient cause or that which is given in reply to the question, ``Where does change (or motion) come from?''. What is singled out in the answer is the whence of change (or motion).
    \item The final cause, the end purpose, is given in reply to the question, ``What is its good?''. What is singled out in the answer is that for the sake of which something is done or takes place.
\end{enumerate}

Aristotle’s framework provided a comprehensive vocabulary for analyzing causality that became foundational to Western scientific and philosophical thought for nearly two millennia.

\subsubsection{Seneca}
\label{sec:history_seneca}

Seneca (c. 56 AD) argues in Letter 65\cite{SenecaLetters} that cause and effect operate within a stage (space) and follow an order (time). Remove the stage or the order, and the conventional understanding of 'making something' or 'causing something' breaks down. His argument highlights time and space as indispensable prerequisites for classical causality. His focus on space and time as necessary conditions served as a precursor for physical concepts that treat spacetime as a background for causal processes.

\subsection{The Realist View: Discovering Causal Laws}
\label{ssec:philosophy_realist}

The realist tradition posits that causal relationships are objective, structural features of the world to be discovered.

\subsubsection{Gottfried Wilhelm Leibniz}
\label{sec:history_leibniz}

The idea of space and time as a background for causality, however, did not remain unchallenged. Gottfried Wilhelm Leibniz (1646--1716) rejected the concept of absolute space and absolute time as independent, fundamental constructs. Instead, Leibniz proposed\cite{LeibnizPhysicsSEP} a relational view in which space is the simultaneous relation of coexisting things and time is the relational order of successive things. Through rigorous first principles analysis, Leibniz argued that the concept of absolute space and time was logically untenable. His relational perspective offered a significant alternative to the preeminent Newtonian worldview of his time.

\subsubsection{John Stuart Mill}
\label{sec:history_stuart_mill}

John Stuart Mill (1806-1873) shifted the focus of causality from metaphysical inquiry to a more practical, methodological problem. In his influential work, \textit{A System of Logic, Ratiocinative and Inductive}\cite{mill2023system}, Mill developed a set of principles, now known as "Mill's Methods," to serve as a guide for discovering causal connections through empirical observation. These methods are designed to systematically eliminate non-causal factors and isolate true causes. Mill's work provided a formal basis for the inductive reasoning that underpins modern experimental science.


\subsubsection{Albert Einstein}
\label{sec:history_einstein}

Albert Einstein (1879--1955) departed from Newtonian physics with his theory of general relativity\cite{EinsteinPapers1915} (GR), in which he established that space and time are one manifold, spacetime, that is bent by the gravitational influence of large masses. General relativity preserves the prerequisite of a spatiotemporal context for causality, echoing Seneca's key insight. However, the notion of a dynamic spacetime requires a dynamic view of causality to fit into the dynamic spacetime manifold.

\subsection{The Empiricist View: Constructing Causal Models}
\label{ssec:philosophy_empiricist}

The empiricist tradition, born from a deep skepticism, argues that causality is not a directly observable force. Instead, it is a mental model that we, as observers, construct to make sense of the statistical regularities we perceive in the world.


\subsubsection{David Hume}
\label{sec:history_hume}

The Scottish philosopher David Hume (1711–1776) launched a more fundamental attack on the notion of causality itself. In his A Treatise of Human Nature\cite{hume2000treatise}, Hume argued that we can never have direct empirical evidence of a necessary connection between a cause and an effect. All we can observe is their constant conjunction in space and time. The idea of "causality" as an unseen force is, for Hume, a product of the human mind developed after observing repeated patterns. Hume's skepticism decouples the idea of causality from any metaphysical necessity and re-grounds it in observational patterns\cite{hume2000treatise}.

\subsubsection{Immanuel Kant}

The German philosopher Immanuel Kant (1724 - 1804) agreed with Hume's premise that causality is not something we can empirically observe in the world. However, he disagreed with Hume's conclusion that it is merely a habit of the mind. In his Critique of Pure Reason \cite{kant2024critique}, Kant argued that causality is an a priori category of understanding. It is a fundamental, built-in "rule" that the mind itself imposes on the raw data of sensory experience to make that experience coherent and understandable. We do not learn causality from the world; rather, we cannot experience the world in an intelligible way without already presupposing the principle of cause and effect. Kant reframes causality from an objective feature of the world (the Realist view) or a mere psychological habit (Hume's view) into a fundamental component of the cognitive architecture of any rational agent. 

\subsubsection{Hans Reichenbach}

Building on Hume's empiricism, Hans Reichenbach (1891 - 1953) was among the first to propose a formal, probabilistic solution to the problem of induction. He argued that causal relationships are not deterministic laws but statistical structures. His 'Principle of the Common Cause' provided the foundational logic for inferring causal structures from statistical correlations\cite{reichenbach1991direction} that later became known as the Reichenbach's Principle used by Pearl to define the concept of a confounder in modern graphical causal models\cite{pearl2000causality}.

\subsubsection{Karl Popper}

The challenge of induction, as framed by Hume, was famously addressed by the philosopher Karl Popper  (1902 – 1994). Popper argued that causal laws can never be verified by observation, only falsified. From this perspective, a causal theory is not a description of a hidden force, but a bold conjecture—a falsifiable model—that is held to be provisionally true only as long as it survives rigorous empirical testing\cite{popper2005logic}.


\subsection{The Practical View: Causality in Physics}
\label{sec:philosophy_physics}

\subsubsection{Bertrand Russell}
\label{sec:history_Russell}

Bertrand Russell (1872--1970) observed that successful physics has its roots in sophisticated, law-based descriptions of how a system evolves dynamically. In modern physics, the focus is on the state of a system (e.g., position, velocity, field strength) and how that entire state evolves continuously and dynamically. Therefore, for Russell, the idea of classical causality, a strict happen-before relation, no longer matches the reality of modern physics. Consequently, Russell wrote in his 1912 essay ``On the Notion of Cause''\cite{RussellOnCause}:  \textit{The law of causality, [...], is a relic of a bygone age.''} 
Many modern physics laws are time-symmetric, which means that if state S1 at time t1 is related to state S2 at time t2 by a law, it is equally true that state S2 at time t2 is related to state S1 at time t1. This relationship is not a simple, linear, one-way street from a necessary ``cause'' to a dependent ``effect.'' Knowing the state at any time allows one to calculate the state at any other time, past or future. Therefore, which state is the ``cause'' and which one is the ``effect'' becomes arbitrary. Russell was not opposed to causality itself; instead, his primary argument was that the traditional philosophical interpretation of causality as a fundamental, temporally asymmetric, and directed link is not what Russell observed in physics. Russell saw physics as a discipline to find, formulate, and test time-symmetric laws of dynamic change. 

\subsubsection{Causality in Quantum Physics}

Russell's view of physics describing time-symmetric laws of dynamic change was taken one step further in quantum physics because there the understanding of causality evolved from a structure that required the pre-existence of space toward a dynamic generative process from which causality emerges\cite{mrini2024indefinitecausalstructurecausal}. This emergent causality does not rely on a pre-existing spacetime but is grounded in a set of underlying rules (i.e., conceptualized as a 'generating function') that determines the fundamental materialization of spatiotemporal properties. The conceptualization of this fundamental level as a ``generating function'' captures the idea of a quantum process from which the necessary condition of classical causality's spatiotemporal structure arises. It is a shift from asking, ``What causes X, given spacetime?'' to ``What process generates spacetime (and thus enables X to be caused)?''.  
 

\subsection{The Process View: Causality as a Process}
\label{sec:philosophy_process}

Aristotle established substance as the primary category of being in his Metaphysics around 350 BC, and every major framework in the Western tradition operated within that metaphysics. Plato's Forms and Aristotle's Four Causes presuppose substances. Seneca treats space and time as the stage on which substances act. Leibniz's monads, despite their relational reframing of space and time, are still defined as simple substances. Mill's methods identify causal connections between events that happen to substances. Hume's skepticism doubts the necessary connection between them, but takes the substances themselves for granted. Kant's a priori categories apply to a world still populated by them. Russell rejects the language of cause but retains a world of states evolving according to law, which is substance metaphysics in a different vocabulary. Even the quantum-causal proposals, which ask what generates spacetime as a background for events, leave the substance metaphysics intact. 

Aristotelian substance has been challenged from two directions in the twentieth century. From the side of language, Ludwig Wittgenstein delivered the most precise and sustained critique of Aristotelian substance. From the side of metaphysics itself, Alfred North Whitehead constructed the first significant alternative. Process philosophy is therefore the first tradition reviewed in this chapter that breaks constructively from the substance metaphysics that prevailed for 2,300 years, a measure of how entrenched the substrate had become.

\subsubsection{Wittgenstein's Critique}

Wittgenstein's critique runs through both phases of his work and converges on a single diagnosis: Aristotelian substance ontology is a projection of grammatical features onto reality, mistaken for an insight into reality. In the Tractatus\cite{wittgenstein1974tractatus}, the early Wittgenstein opens with the claim that the world is "the totality of facts, not of things" (1.1), inverting the Aristotelian picture in a single sentence. Aristotelian primary substance is supposed to exist in itself and to bear properties; Wittgenstein's objects, by contrast, have no in-itself and are exhausted by their possibilities of combination with other objects in states of affairs (2.0123). The unit of reality is the obtaining of a fact. In this view, reading substance metaphysics off grammatical structure is a category error.
The late Wittgenstein, in the Philosophical Investigations\cite{wittgenstein2009investigations}, intensifies the critique from the side of meaning. Meaning is constituted by use within a form of life, not by correspondence to hidden essential structure. Once this is seen, the Aristotelian apparatus of substances bearing essences to ground the meanings of general terms becomes unnecessary scaffolding. In his view, substance metaphysics is what one gets by mistaking how we talk for how things are.

\subsubsection{Whitehead's Process and Reality}

The process view, originating in Whitehead's Process and Reality\cite{whitehead2010process} and developed further by Bergson\cite{bergson2022creative}, rejects the Aristotelian substance metaphysics. For the process philosopher, the basic constituents of reality are not substances that undergo change but events that \emph{are} their change. On this view, causality is not a relation linking two pre-existing things. It is internal to the constitution of an event. Whitehead's actual occasions are momentary, atomic processes of becoming, each of which inherits from prior occasions and is then taken up as data by later ones. Cause and effect are not two separate categories but two roles successively occupied by the same occasion. The category of substance is replaced by what Whitehead calls the Category of the Ultimate: creativity, the many, and the one\cite{whitehead2010process}. Whitehead proposes the first comprehensive non-Aristotelian metaphysical system in Western philosophy. The proposal is comprehensive in scope. Substance ontology underwrites the philosophy of causality, the subject-predicate form of expression, the logic of attribution, the epistemology of objects, and the standard picture of scientific law as describing how things behave. Whitehead challenges all of it at once, which is part of why Process and Reality is famously difficult. It is also the immediate philosophical ancestor of the Functional View.
  

\subsection{The Functional View of Causality}
\label{sec:philosophy_functional}

Russell correctly observed that classical causality, with its assumption
of temporal precedence as a constitutive feature of the cause–effect
relation, does not fit modern physics, which models dynamic change using
time-symmetric equations. The conflict is that classical causality
requires that the cause precede the effect. However, the laws of motion in
classical mechanics, the Schrödinger equation in quantum mechanics, and
Maxwell's equations in electromagnetism are all time-symmetric, meaning
they relate states without distinguishing past from future. 

The asymmetry that classical causality treats as fundamental was articulated clearly by Seneca, who argued that cause and effect operate within a stage (space) and follow an order (time) in which substances act on substances. Whitehead's atomic processes of becoming, by contrast, inherit from prior occasions and are taken up as data by later ones, without requiring a fixed spatiotemporal stage.


The author agrees with the premise established by Russell and embraces the process view established by Whitehead. The functional view of causality generalizes causality as a spacetime agnostic functional dependency codified as the effect propagation process, the EPP. The EPP establishes that  effects propagate asymmetrically in sequence, and that the laws governing 
the changes within an effect do not need to carry any assumption about time symmetry. The propagation of effects is constituted by the composition of functions: each step of a causal process is a function from one propagating effect to the next one, and a causal model is the composition of those steps into a process. Monadic composition preserves the causal order of effects throughout the process. In that sense, the functional view encodes process causality.    
 
\subsubsection{The Axiom of Process Causality}

The EPP is built on a single axiom from which the entire framework is
derived. The axiom states that one propagating effect is obtained from
another by monadic composition encapsulating a causal function:
 
\begin{quotation}
\textbf{$m_{2} = m_{1} \mathbin{>\!\!>\!\!=} f$}
\end{quotation}


\noindent Here $m_1$ and $m_2$ are propagating effects, $f$ is a causal
function from a value to a propagating effect, and $\mathbin{>\!\!>\!\!=}$
(read ``bind'') is the monadic composition operator. A propagating effect
is a structured object that carries, alongside its value, the state, the
context, the error condition, and the log of the causal process up
to that point. The bind operation carries all of these through the
composition as is standard in monadic composition. The detailed formulation of bind, the structure of the propagating effect, 
and the higher-kinded type machinery that permits the same monad to operate over
deterministic, probabilistic, and time-symmetric functions
are developed in section \ref{sec:epp} 

\subsubsection{Consequences of the Single Axiom}

The single monadic axiom yields a cascade of consequences from which much of the
EPP follows:

\begin{enumerate}
	\item \textbf{Inherently Agnostic}: The causal function $f$ and the
	values carried by propagating effects are unconstrained. A function
	may be deterministic, probabilistic, symbolic, or numerically
	intensive; the monad imposes no further restriction.
	
	\item \textbf{Inherently Computable}: Functions compute, and monadic
	composition is a computable operation. The axiom is therefore
	directly programmable; what the axiom states, the implementation
	executes.
	
	\item \textbf{Inherently Composable}: Monadic composition is the most
	general form of functional composition that carries context with it.
	Where ordinary function composition gives $g \circ f$, monadic
	composition through bind gives the same composition with state,
	error, and log threaded automatically. Category theory, which
	supplies the formal foundations of monadic composition, is well
	established.
	
	\item \textbf{Inherently Dynamic}: The state component of the
	propagating effect carries information that evolves through the
	composition chain. A causal model expressed in the EPP is therefore
	dynamic by default rather than dynamic by extension.
	
	\item \textbf{Inherently Asymmetric in Propagation}: The bind operator
	is unidirectional by construction. Given a chain
	$m \mathbin{>\!\!>\!\!=} f \mathbin{>\!\!>\!\!=} g$, the composition
	moves from left to right and admits no inverse. This supplies the
	propagation asymmetry that the classical definition of causality
	requires, without imposing any constraint on the laws inside each
	step. The function $f$ may be a time-symmetric physics equation, a
	discrete logical implication, or a stochastic transition; the
	directional flow is carried by the monadic structure rather than by
	an assumption about the laws being propagated.
	
	\item \textbf{Inherently Subsumptive}: The EPP rests
on a single axiom and therefore subsumes any causal framework built on
a stronger axiomatic basis. Pearl's structural causal models, Reichenbach's
common-cause principle, the Suppes probabilistic framework, and the
classical definition itself are recovered from the EPP by adjoining the
additional axioms that distinguish each. 

\item \textbf{Inherently Emergent}: Structural properties that other
causal frameworks impose as foundational axioms appear in the EPP only
as properties of specific causal models. Acyclicity, which Pearl's directed acyclic graphs require by construction, holds in the EPP for those causal models that are acyclic
and is absent for those that contain feedback. Time symmetry holds for
specific functions (e.g Maxwell's equations) and not
for others (thermodynamic descriptions, classical implication).
Determinism holds for deterministic functions and not for probabilistic
or stochastic ones. The Markov property holds for monadic contexts in
which the state field carries the relevant information and not
otherwise. The framework imposes none of these; they emerge, when they
emerge.  This is the formal basis for the metaphysics of emergence developed in \ref{sec:metaphysics_becoming}
\end{enumerate}

The fifth consequence resolves what would otherwise be a difficult
reconciliation of the EPP with Russell's critique. Russell objected that
classical causality assumes temporal asymmetry, while modern physics uses
time-symmetric equations. Russell is right, and so is the classical definition of causality.  

However, the underlying problem rests on two distinct asymmetries. The first is the asymmetry of laws: whether the equations governing a system are invariant under time reversal. The second is the asymmetry of propagation: the directional flow by which the output of one causal step becomes the input of the next. Russell's critique applies to the first based on the physics perspective; Seneca's, Aristotle's, and Pearl's classical definitions of causality require the second. These two asymmetries are independent from each other and the  monadic axiom separates them via its bind operator. The bind operator treats propagation asymmetry as structural integrity, while the functions inside the bind remain free of any assumption of time. The EPP
therefore preserves the directional order of classical causality (the flow from cause to effect) while also preserving the time-symmetric equations of modern physics. 
 
\subsection{Discussion}
\label{sec:philosophy_discussion} 

The Functional View advances the philosophy of causality on two specific points where the prior traditions surveyed in this chapter share a common limitation. The first concerns the structure of causal asymmetry. The classical, empiricist, and realist traditions all treat causal asymmetry as a single property: cause precedes effect, the laws describing the relation respect that order, and the relation has a direction because that's how it was since Aristotle. 

Russell's critique exploits this conflation. He observes that the laws of modern physics are time-symmetric and concludes that classical causality, with its assumption of asymmetry, is a relic of a bygone era.
 
The Functional View shows the conflation is unnecessary. Two asymmetries are at work, and they are independent. The first is the asymmetry of laws, which is what time-symmetric equations lack. The second is the asymmetry of propagation, which is the directional flow from cause to effect. Russell's critique applies to the first; Seneca's, Aristotle's, and Pearl's classical definitions require the second. Once the two are separated, the conflict between classical causality and modern physics resolves. The Functional View carries propagation asymmetry in the composition operator and leaves the laws inside each step unconstrained.
  
Seneca's argument that cause and effect operate within a stage and follow an order has shaped every realist and empiricist position ever since. Even Russell, who rejects the language of cause, retains a world of states evolving according to law within a spatiotemporal background. The process tradition is the first to question whether such a background is required, but Whitehead's actual occasions and Bergson's \emph{durée} indicate the dispensability of the stage rather than demonstrate it. The Functional View completes the move. By treating space and time as one possible kind of context that a propagating effect can carry, it shows how a causal process can run without a presupposed spatiotemporal manifold. Spacetime becomes a feature of certain causal models rather than a precondition for causality as such. This is the operational form of the metaphysical commitment Whitehead names.

The functional view does not resolve the metaphysical disagreement among the realist, empiricist, and process traditions on what causality fundamentally is. A realist may read the function inside a bind step as an objective structural feature of the world. An empiricist may read it as the formal representation of a model inferred from data. A process philosopher may read it as the determinate content of an actual occasion taking up its predecessors. The Functional View accommodates all three readings out of the recognition that the asymmetry-separation and the spacetime-detachment moves can be made without committing to a particular metaphysics of what the function represents.  

\newpage